\chapter{Całki funkcji wektorowych}
\label{ma2:ch:vector-integrals}

W rozdziale \ref{ma2:ch:limits-calculus}
wprowadziliśmy całkę Riemanna krzywych o wartościach
w przestrzeniach unormowanych, a w rozdziale
\ref{ma2:ch:primitives} wykazaliśmy addytywność,
kryterium Cauchy'ego i wzór Newtona--Leibniza.
Uzupełnimy teraz teorię o własności charakterystyczne
dla funkcji wektorowych, w szczególności o różnice
między przypadkiem skończonego i nieskończonego
wymiaru.

Ogólna teoria całki Riemanna w przestrzeniach
lokalnie wypukłych znajduje się w
\atlasref{FA-15A::sec:integral-lcs15a}{analizie funkcjonalnej, moduł 15A}.
Całki Bochnera, Birkhoffa, Pettisa, Dunforda
i Gelfanda omówiono w
\atlasref{FA-15B::sec:integral-vector15b}{module 15B},
a teorię miar wektorowych i własność
Radona--Nikodyma również w
\atlasref{ME-30::ch:vector-measures-30}{teorii miary i całki, rozdział 30}.
Nie będziemy powielać tych konstrukcji i ich
ogólnej teorii miarowej.

\section{Całka Riemanna w przestrzeni skończenie wymiarowej}

Niech $a<b$, $E$ będzie przestrzenią Banacha
nad $\R$ lub $\C$, a $f:[a,b]\to E$.
Sumy Riemanna oznaczamy jak w rozdziale
\ref{ma2:ch:primitives} przez $R(f,P)$,
gdzie $P$ jest podziałem z punktami pośrednimi,
a $|P|$ oznacza jego średnicę.
Wszystkie całki Riemanna rozpatrujemy w normie $E$.

\begin{twierdzenie}[Całkowanie współrzędnych]
\label{ma2:thm:vector-integral-coordinates}
Niech $E$ będzie rzeczywistą przestrzenią
skończenie wymiarową z bazą $(e_1,\ldots,e_d)$.
Dla $f(t)=\sum_{j=1}^d f_j(t)e_j$
całkowalność Riemanna $f$ jest równoważna
całkowalności każdej współrzędnej $f_j$.
Wówczas
\[
 \int_a^b f(t)\,\mathrm dt
 =\sum_{j=1}^d e_j\int_a^b f_j(t)\,\mathrm dt.
\]
Dla zespolonej przestrzeni skończenie wymiarowej
można zastosować ten wynik do współrzędnych
rzeczywistych i urojonych.
\end{twierdzenie}
\begin{proof}
Ciągłość funkcjonałów współrzędnych $\pi_j:E\to\R$
wynika z równoważności norm w skończonym wymiarze.
Jeżeli $f$ jest całkowalna, to $f_j=\pi_j\circ f$
jest całkowalna, a jej całka jest $j$-tą
współrzędną całki $f$, z przemienności całki
z operatorem liniowym, wykazanej po twierdzeniu
\ref{ma2:thm:primitive-subinterval}.

Odwrotnie, niech $I_j$ będą całkami współrzędnych.
Z liniowości skończonych sum
\[
 \left\|R(f,P)-\sum_{j=1}^d I_je_j\right\|
 \le\sum_{j=1}^d\|e_j\|\,|R(f_j,P)-I_j|.
\]
Dla $\varepsilon>0$ i $d\ge1$ wybieramy
$\delta_j>0$ tak, aby
$|R(f_j,P)-I_j|<\varepsilon/(d\|e_j\|)$
dla $|P|<\delta_j$. Dla
$|P|<\min_j\delta_j$ prawa strona jest
mniejsza od $\varepsilon$.
Jeżeli $d=0$, wszystkie sumy i całki są zerowe.
W przypadku zespolonym traktujemy $E$
jako rzeczywistą przestrzeń skończenie wymiarową;
norma i sumy Riemanna pozostają te same.
\end{proof}

\begin{wniosek}[Skończenie wymiarowe kryterium Lebesgue'a]
\label{ma2:cor:vector-integral-lebesgue-finite}
Niech $E$ będzie skończenie wymiarowa,
a $f:[a,b]\to E$ ograniczona.
Wówczas $f$ jest całkowalna w sensie Riemanna
wtedy i tylko wtedy, gdy jest ciągła w normie
poza zbiorem miary Lebesgue'a zero.
\end{wniosek}
\begin{proof}
Z ciągłości funkcjonałów współrzędnych
ciągłość $f$ w danym punkcie pociąga
ciągłość każdej współrzędnej. Odwrotnie,
ciągłość wszystkich współrzędnych daje
ciągłość $f$ na mocy oszacowania
$\|f(t)-f(s)\|\le\sum_j\|e_j\|\,|f_j(t)-f_j(s)|$.
Zatem zbiór punktów nieciągłości $f$
jest skończoną sumą zbiorów punktów
nieciągłości współrzędnych.
Stosujemy kryterium Lebesgue'a dla funkcji
skalarnych (twierdzenie
\ref{ma2:thm:riemann-box-lebesgue-criterion}
dla $n=1$) oraz twierdzenie
\ref{ma2:thm:vector-integral-coordinates}.
Skończona suma zbiorów zerowych jest zerowa,
a podzbiór zbioru zerowego jest mierzalny
i zerowy względem zupełnej miary Lebesgue'a.
\end{proof}

\begin{przyklad}[Krzywa o wartościach w $\R^3$]
\label{ma2:ex:vector-integral-coordinates}
Dla $f(t)=(\cos t,\sin t,t)$ i $0\le t\le2\pi$
każda współrzędna jest ciągła. Zatem
\[
 \int_0^{2\pi} f(t)\,\mathrm dt
 =(0,0,2\pi^2).
\]
Istotnie, funkcjami pierwotnymi współrzędnych
są $\sin t$, $-\cos t$ i $t^2/2$.
Pierwsze dwie składowe znoszą się po pełnym
obrocie, lecz trzecia ma dodatnią całkę.
\end{przyklad}

\section{Oszacowania normy i całkowalność w nieskończonym wymiarze}

\begin{stwierdzenie}
\label{ma2:prop:vector-integral-norm}
Jeżeli zarówno $f:[a,b]\to E$, jak i skalarna
funkcja $t\mapsto\|f(t)\|$ są całkowalne
w sensie Riemanna, to
\[
 \left\|\int_a^b f(t)\,\mathrm dt\right\|
 \le\int_a^b\|f(t)\|\,\mathrm dt.
\]
W szczególności obie całki istnieją dla
ciągłej $f$.
\end{stwierdzenie}
\begin{proof}
Nierówność trójkąta daje dla każdego
podziału z punktami pośrednimi $P$
\[
 \|R(f,P)\|
 \le\sum_j(t_j-t_{j-1})\|f(\xi_j)\|
 =R(\|f\|,P).
\]
Biorąc dowolny ciąg podziałów o średnicach
dążących do zera, przechodzimy do granicy
po obu stronach, korzystając z ciągłości normy.
Jeżeli $f$ jest ciągła, to funkcja $\|f\|$
jest ciągła z nierówności
$|\|x\|-\|y\||\le\|x-y\|$.
Istnienie całki Riemanna funkcji ciągłej
o wartościach w przestrzeni Banacha wynika z
\atlasref{FA-15A::thm:riemann-existence15a}{twierdzenia o istnieniu całki}.
\end{proof}

W skończonym wymiarze całkowalność $f$
zapewnia automatycznie całkowalność $\|f\|$:
funkcja normy jest ograniczona i ciągła
w każdym punkcie ciągłości $f$,
więc wystarczy zastosować wniosek
\ref{ma2:cor:vector-integral-lebesgue-finite}
i skalarne kryterium Lebesgue'a.
W nieskończonym wymiarze dodatkowa
hipoteza o całkowalności normy nie jest
zbędna.

\begin{przyklad}[Całkowalna funkcja o niecałkowalnej normie]
\label{ma2:ex:vector-integral-norm-failure}
Niech $E=\ell^2(\mathbb N)$ z bazą
ortonormalną $(e_k)_{k\ge1}$.
Ustalmy bijektywną enumerację $(q_k)_{k\ge1}$
zbioru $\Q\cap[0,1]$ i połóżmy
\[
 f(t)=
 \begin{cases}
 e_k,&t=q_k,\\
 0,&t\notin\Q.
 \end{cases}
\]
Funkcja ta jest całkowalna w sensie Riemanna
i ma całkę zero, mimo że jest nieciągła
w normie w każdym punkcie, a
$\|f(t)\|=\chi_{\Q\cap[0,1]}(t)$
nie jest całkowalna w sensie Riemanna.

Aby wykazać całkowalność, dla dowolnego
podziału z punktami pośrednimi oznaczmy
przez $\alpha_k$ sumę długości podprzedziałów,
których punkt pośredni jest równy $q_k$.
Tylko skończenie wiele $\alpha_k$ jest
niezerowych i
$R(f,P)=\sum_k\alpha_ke_k$.
Każdy punkt może należeć do najwyżej
dwóch podprzedziałów podziału, więc
$0\le\alpha_k\le2|P|$. Ponadto
$\sum_k\alpha_k\le1$, gdyż podprzedziały
ze znacznikami niewymiernymi nie wnoszą wkładu.
Z ortonormalności mamy
\[
 \|R(f,P)\|^2=\sum_k\alpha_k^2
 \le 2|P|\sum_k\alpha_k\le2|P|.
\]
Wszystkie sumy dążą do zera wraz
ze średnicą podziału, jednostajnie
względem wyboru punktów pośrednich.

Każde względne otoczenie punktu $[0,1]$
zawiera punkty wymierne i niewymierne.
Normy wartości wynoszą tam odpowiednio $1$ i $0$,
więc norma nie jest ciągła w żadnym punkcie.
Ciągłość $f$ pociągałaby ciągłość normy,
zatem także $f$ nie jest nigdzie ciągła.
Dolne sumy Darboux funkcji $\|f\|$ są
zerowe, górne równe $1$, dlatego norma
nie jest całkowalna w sensie Riemanna.
\end{przyklad}

Przykład \ref{ma2:ex:vector-integral-norm-failure}
pokazuje również, że konieczność wniosku
\ref{ma2:cor:vector-integral-lebesgue-finite}
nie przechodzi automatycznie na przestrzenie
Banacha nieskończonego wymiaru, nawet
na ośrodkową przestrzeń Hilberta.
Dostateczność natomiast pozostaje prawdziwa.

\begin{twierdzenie}[Warunek dostateczny całkowalności]
\label{ma2:thm:vector-riemann-sufficient}
Niech $f:[a,b]\to E$ będzie ograniczona
i ciągła w normie poza zbiorem miary
Lebesgue'a zero. Wtedy $f$ jest całkowalna
w sensie Riemanna.
\end{twierdzenie}
\begin{proof}
Połóżmy $L=b-a$ i wybierzmy $M>0$
takie, że $\|f(t)\|\le M$.
Dla $x\in[a,b]$ określamy oscylację punktową
\[
 \omega_f(x)=\inf_{r>0}\sup\{
 \|f(s)-f(t)\|:s,t\in[a,b]\cap(x-r,x+r)\}.
\]
Z definicji ciągłość $f$ w $x$ jest
równoważna warunkowi $\omega_f(x)=0$:
w jedną stronę stosujemy nierówność
trójkąta przez $f(x)$, w drugą bierzemy
jeden z argumentów równy $x$.
Dla $\eta>0$ zbiór
$D_\eta=\{x:\omega_f(x)\ge\eta\}$
jest domknięty. Jeżeli $\omega_f(x)<\eta$,
to istnieje otoczenie $x$ o oscylacji
mniejszej od $\eta$; każde $y$ z tego
otoczenia ma mniejsze otoczenie w nim
zawarte, więc $\omega_f(y)<\eta$.
Zatem $D_\eta$ jest zwarty i ma miarę zero.

Ustalmy $\varepsilon>0$ i wybierzmy
$\eta=\varepsilon/(4L)$.
Z regularności zewnętrznej miary Lebesgue'a
istnieje otwarty $O\subset\R$ zawierający
$D_\eta$ o mierze mniejszej od $\varepsilon/(8M)$.
Poza $D_\eta$ każdy punkt ma względnie
otwarte otoczenie o oscylacji mniejszej
od $\eta$. Otoczenia te wraz z
$O\cap[a,b]$ pokrywają zwarty przedział.
Pokrycie ma skończone podpokrycie i liczbę
Lebesgue'a $\delta>0$ (por. dowód
twierdzenia \ref{ma2:thm:riemann-box-lebesgue-criterion}).

Niech $P,T$ będą dowolnymi podziałami
o średnicach mniejszych od $\delta$.
Na wspólnym rozdrobnieniu wybierzmy dowolne
punkty pośrednie i oznaczmy otrzymany
podział przez $S$. W różnicy $R(f,P)-R(f,S)$
każda wartość przypisana podprzedziałowi $P$
jest porównywana z wartościami na jego
podprzedziałach w $S$. Jeżeli podprzedział $P$
zawiera się w $O$, norma każdej różnicy wartości jest
ograniczona przez $2M$.
W przeciwnym razie podprzedział zawiera się
w jednym z pozostałych elementów podpokrycia,
więc każda norma różnicy jest mniejsza od $\eta$.
Wkłady ważą się długościami podprzedziałów,
a łączna długość podprzedziałów $P$ zawartych
w $O$ nie przekracza $\lambda_1(O)$, gdyż
ich wnętrza są rozłączne.
Zatem
\[
 \|R(f,P)-R(f,S)\|
 \le 2M\lambda_1(O)+\eta L<\varepsilon/2.
\]
Analogiczne oszacowanie dla $T$ daje
$\|R(f,T)-R(f,S)\|<\varepsilon/2$.
Z nierówności trójkąta otrzymujemy
$\|R(f,P)-R(f,T)\|<\varepsilon$.
Kryterium Cauchy'ego
\ref{ma2:lem:primitive-cauchy} daje tezę.
\end{proof}

% Następny blok: porównanie całek Riemanna i Bochnera,
% całki słabe i silne, przykłady funkcyjne i operatorowe.
